\chapter{}
\label{chap:3}

\section{Умова}
Знайти найбільш правдоподібний розв'язок системи лінійних рівнянь зі спотвореними правими частинами $Ax = b$, якщо:
\begin{equation*}
    A = \begin{pmatrix}
        0 & 0 & 1 \\
        1 & 1 & 1 \\
        0 & 1 & 1 \\
        1 & 1 & 0 \\
        1 & 0 & 0
    \end{pmatrix}, \quad
    b = \begin{pmatrix}
        0 \\ 1 \\ 0 \\ 1 \\ 0
    \end{pmatrix}
\end{equation*}

\section{Розв'язання}
Для знаходження розв'язку будуємо емпіричний розподіл $g(y)$ за формулою:
\begin{equation*}
    g(y) = \sum\limits_{\substack{i = \overline{1, m} \\ A_i = y}} (-1)^{b_i}
\end{equation*}
де $A_i$ -- $i$-й рядок матриці $A$, $b_{i} \in \left\{0, 1\right\}$, $y \in V_{n}$.

\noindent Побудуємо таблицю істинності і функцію $g(y)$
\begin{table}[htpb]
    \centering
    \begin{tblr}{
        colspec = {c | c | l | c | c},
        row{1} = {font=\bfseries},
        hline{1,2,10} = {solid}
        }
        $y$   & Крок 1 & Збіги                     & Доданок                      & $g(y)$ \\
        $000$ & 0      & Немає збігів з матрицею A & $+0$                         & $0$    \\
        $001$ & 0      & Збіг у 1 рядку            & $+(-1)^{b_1} = (-1)^0 = 1$   & $1$    \\
        $010$ & 0      & Немає збігів              & $+0$                         & $0$    \\
        $011$ & 0      & Збіг у 3 рядку            & $+ (-1)^{b_3} = (-1)^0 = 1$  & $1$    \\
        $100$ & 0      & Збіг у 5 рядку            & $+ (-1)^{b_5} = (-1)^0 = 1$  & $1$    \\
        $101$ & 0      & Немає збігів              & $+0$                         & $0$    \\
        $110$ & 0      & Збіг у 4 рядку            & $+ (-1)^{b_4} = (-1)^1 = -1$ & $-1$   \\
        $111$ & 0      & Збіг у 2 рядку            & $+ (-1)^{b_2} = (-1)^1 = -1$ & $-1$   \\
    \end{tblr}
    \caption{Еміричний розподіл вектора $g(y)$}
\end{table}

Тепер застосуємо швидке перетворення Уолша-Адамара, яке тренували у минулому дз, до вектора $g$, щоб отримати його 
коефіцієнти Фур'є $C_{g}$. Найбільші значення у спектрі вкажуть на найбільш правдоподібні розв'язки.

\begin{table}[htpb]
    \centering
    \begin{tblr}{
        colspec = {c | c | c | c},
        row{1} = {font=\bfseries},
        hline{2,10} = {solid}, % pivot lines
        hline{6} = {1}{solid}, % first column
        hline{4,6,8} = {2}{solid}, % second column
        hline{2,3,4,5,6,7,8,9} = {3}{solid} % third column
            }
        $g$  & Крок 1 & Крок 2 & $C_{g}$        \\
        0    & 1      & 0      & 1              \\
        1    & 1      & 1      & $-1$           \\
        0    & $-1$   & 2      & 3 \\
        1    & 0      & 1      & $-1$           \\
        1    & $-1$   & 0      & 3 \\
        0    & 1      & 3      & $-3$           \\
        $-1$ & 1      & $-2$   & $-3$           \\
        $-1$ & 2      & $-1$   & $-1$           \\
    \end{tblr}
    \caption{Швидке перетворення Адамара для $f = (0,1,0,1,1,0,-1,-1)$}
\end{table}
На всяк випадок зроблю перевірку рівності Парсеваля. У нас $\operatorname{wt}(g) = 5$ -- сума квадратів значень 
функції $g$.
\begin{equation*}
    \frac{1^2 + (-1)^2 + 3^2 + (-1)^2 + 3^2 + (-3)^2 + (-3)^2 + (-1)^2}{2^3}
    = \frac{40}{8} = 5 = \operatorname{wt}(g).
\end{equation*}
Отже коефіцієнти Фур'є нашої функції це: $C_{g} = \left(1, -1, 3, -1, 3, -3, -3, -1\right)$. Максимальне значення 
у рядку це $3$, і воно досягається на двох позиціях -- третя і п'ята. Тобто
\begin{itemize}
    \item $C_{g}(010) = 3 \implies y_{1} = (0, 1, 0)$
    \item $C_{g}(100) = 3 \implies y_{2} = (1, 0, 0)$
\end{itemize}

Отже, найбільш правдоподібними розв'язками системи є вектори $(0, 1, 0)$ та $(1, 0, 0)$.
